% ===== PRÉAMBULE CANONIQUE — à coller VERBATIM en tête de CHAQUE .tex =====
\documentclass[11pt,a4paper]{article}
\usepackage[utf8]{inputenc}\usepackage[T1]{fontenc}\usepackage{lmodern}
\usepackage[french]{babel}
\usepackage{amsmath,amssymb,mathtools}\usepackage{icomma}
\usepackage{siunitx}\sisetup{locale=FR}  % \qty{}{}, \si{}, \ang{}
\usepackage{xcolor}\usepackage{colortbl}
\usepackage{tikz}\usepackage{pgfplots}\pgfplotsset{compat=1.18}
\usepackage[siunitx,europeanresistors]{circuitikz} % circuits (élec.)
\usepackage[most]{tcolorbox}\usepackage{enumitem}\usepackage{array,booktabs}
\usepackage[a4paper,margin=2.2cm]{geometry}
\usetikzlibrary{arrows.meta,calc,angles,quotes,positioning,patterns,%
  backgrounds,shapes.geometric,decorations.pathmorphing,decorations.markings,intersections}
\definecolor{primary}{RGB}{222,49,99}\definecolor{primarydark}{RGB}{150,22,55}
\definecolor{defcol}{RGB}{0,114,178}\definecolor{methcol}{RGB}{52,92,148}
\definecolor{excol}{RGB}{0,135,90}\definecolor{attncol}{RGB}{200,40,55}
\tcbset{boxbase/.style={enhanced,breakable,boxrule=0.9pt,arc=2.5pt,
  left=3mm,right=3mm,top=2.3mm,bottom=2.3mm,fonttitle=\bfseries,coltitle=white}}
% --- boîtes TUTORIEL (noms EXACTS, ne pas renommer) ---
\newtcolorbox{cle}[1][]{boxbase,colback=defcol!6,colframe=defcol,
  title={\textsc{À retenir}\ifx\relax#1\relax\relax\else\textmd{ : \itshape #1}\fi}}
\newtcolorbox{methode}[1][]{boxbase,colback=methcol!7,colframe=methcol,sharp corners,
  title={\textsc{Méthode}\ifx\relax#1\relax\relax\else\textmd{ --- \itshape #1}\fi}}
\newtcolorbox{exemplebox}[1][]{boxbase,colback=excol!5,colframe=excol,
  title={\textsc{Exemple}\ifx\relax#1\relax\relax\else\textmd{ : \itshape #1}\fi}}
\newtcolorbox{piege}[1][]{boxbase,colback=attncol!6,colframe=attncol,
  title={\textsc{Piège}\ifx\relax#1\relax\relax\else\textmd{ : \itshape #1}\fi}}
\newtcolorbox{remarque}[1][]{boxbase,colback=black!4,colframe=black!45,
  title={\textsc{Remarque}\ifx\relax#1\relax\relax\else\textmd{ : \itshape #1}\fi}}
% --- boîtes EXERCICE / CORRIGÉ (noms EXACTS, auto-comptées) ---
\newtcolorbox[auto counter]{exo}[1][]{boxbase,colback=primary!4,colframe=primary,
  title={\textsc{Exercice}~\thetcbcounter\ifx\relax#1\relax\relax\else\textmd{ --- \itshape #1}\fi}}
\newtcolorbox[auto counter]{corrige}[1][]{boxbase,colback=excol!4,colframe=excol,
  title={\textsc{Corrigé}~\thetcbcounter\ifx\relax#1\relax\relax\else\textmd{ --- \itshape #1}\fi}}
\newcommand{\vv}[1]{\vec{#1}}\newcommand{\ds}{\displaystyle}
\newcommand{\R}{\mathbb{R}}\newcommand{\N}{\mathbb{N}}\newcommand{\Z}{\mathbb{Z}}
% (un \usepackage{fancyhdr}+en-tête est possible ; il est ignoré côté web)
% =========================================================================
\begin{document}
\sisetup{per-mode=symbol}

\begin{exo}[quel est le signe de la charge ?]
Voici deux spectres de charges ponctuelles isolées.
\begin{center}
\begin{tikzpicture}[scale=0.62,>=Stealth,font=\scriptsize,
    ch/.style={circle,minimum size=6.5mm,inner sep=0pt,text=white,font=\bfseries\scriptsize}]
  \begin{scope}[xshift=-3.2cm]
    \foreach \a in {0,45,...,315}{\draw[->,excol,line width=0.8pt] (0,0)--(\a:1.7);}
    \node[ch,fill=black!45] at (0,0) {?};
    \node[primarydark,font=\bfseries] at (0,-2.4) {spectre I};
  \end{scope}
  \begin{scope}[xshift=3.2cm]
    \foreach \a in {0,45,...,315}{\draw[->,excol,line width=0.8pt] (\a:1.7)--(0,0);}
    \node[ch,fill=black!45] at (0,0) {?};
    \node[primarydark,font=\bfseries] at (0,-2.4) {spectre II};
  \end{scope}
\end{tikzpicture}
\end{center}
\begin{enumerate}[label=\alph*)]
  \item Quel est le signe de la charge du spectre I ?
  \item Quel est le signe de la charge du spectre II ?
\end{enumerate}
\end{exo}

\begin{exo}[lecture d'un dipôle]
Le spectre ci-dessous est celui de deux charges opposées.
\begin{center}
\begin{tikzpicture}[scale=0.8,>=Stealth,font=\scriptsize,
    ch/.style={circle,minimum size=7mm,inner sep=0pt,text=white,font=\bfseries\scriptsize}]
  \node[ch,fill=black!45] (p) at (-1.6,0) {$X$};
  \node[ch,fill=black!45] (n) at (1.6,0) {$Y$};
  \draw[->,excol,line width=0.8pt] (-1.25,0)--(1.25,0);
  \foreach \s in {0.75,1.4}{
    \draw[->,excol,line width=0.8pt] (-1.35,0.16*\s) .. controls (-0.6,\s) and (0.6,\s) .. (1.35,0.16*\s);
    \draw[->,excol,line width=0.8pt] (-1.35,-0.16*\s) .. controls (-0.6,-\s) and (0.6,-\s) .. (1.35,-0.16*\s);}
\end{tikzpicture}
\end{center}
\begin{enumerate}[label=\alph*)]
  \item Laquelle des deux charges ($X$ ou $Y$) est positive ?
  \item Justifie en une phrase à partir du sens des lignes.
\end{enumerate}
\end{exo}

\begin{exo}[où le champ est-il le plus fort ?]
Autour d'une charge ponctuelle, on repère deux points : $P$ (proche de la charge) et $Q$ (éloigné).
\begin{center}
\begin{tikzpicture}[scale=0.7,>=Stealth,font=\scriptsize,
    ch/.style={circle,minimum size=6.5mm,inner sep=0pt,text=white,font=\bfseries\scriptsize}]
  \foreach \a in {0,45,...,315}{\draw[->,excol,line width=0.8pt] (0,0)--(\a:2.4);}
  \node[ch,fill=attncol] at (0,0) {$+$};
  \node[circle,fill=black,inner sep=1.3pt,label=above:{$P$}] at (0.95,0.0) {};
  \node[circle,fill=black,inner sep=1.3pt,label=above:{$Q$}] at (2.15,0.0) {};
\end{tikzpicture}
\end{center}
\begin{enumerate}[label=\alph*)]
  \item En quel point le champ est-il le plus intense ? Justifie par la densité des lignes.
  \item Si l'on \emph{double} la charge, les lignes de champ deviennent-elles plus denses ou plus
        espacées ?
\end{enumerate}
\end{exo}

\begin{exo}[le condensateur plan]
Un condensateur plan a son armature \textbf{supérieure} chargée $+$ et son armature
\textbf{inférieure} chargée $-$.
\begin{center}
\begin{tikzpicture}[scale=0.9,>=Stealth,font=\scriptsize]
  \fill[attncol!35] (-2.4,1.1) rectangle (2.4,1.25); \node[attncol,font=\bfseries] at (2.9,1.18){$+$};
  \fill[defcol!35] (-2.4,-1.25) rectangle (2.4,-1.1); \node[defcol,font=\bfseries] at (2.9,-1.18){$-$};
  \foreach \x in {-1.8,-0.9,0,0.9,1.8}{\draw[->,excol,line width=0.9pt] (\x,1.0)--(\x,-1.0);}
  \node[left,font=\scriptsize] at (-2.5,1.18){$A$}; \node[left,font=\scriptsize] at (-2.5,-1.18){$B$};
\end{tikzpicture}
\end{center}
\begin{enumerate}[label=\alph*)]
  \item Dans quel sens est orienté le champ $\vv{E}$ entre les armatures ?
  \item Quelle armature ($A$ ou $B$) est au potentiel le plus élevé ?
\end{enumerate}
\end{exo}

\begin{exo}[une propriété interdite]
On propose le spectre suivant, où deux lignes de champ \textbf{se croisent} au point $X$.
\begin{center}
\begin{tikzpicture}[scale=0.9,>=Stealth,font=\scriptsize]
  \draw[->,excol,line width=0.9pt] (-2,-1) .. controls (-0.3,0.2) and (0.3,0.2) .. (2,-1);
  \draw[->,excol,line width=0.9pt] (-2,1) .. controls (-0.3,-0.2) and (0.3,-0.2) .. (2,1);
  \node[circle,fill=black,inner sep=1.3pt,label=below:{$X$}] at (0,0.02) {};
\end{tikzpicture}
\end{center}
\begin{enumerate}[label=\alph*)]
  \item Ce spectre est-il physiquement possible ? Justifie.
  \item Quelle propriété fondamentale des lignes de champ est violée en $X$ ?
\end{enumerate}
\end{exo}

\end{document}
